Nationwide internet shutdowns are a recurring instrument of state control,
documented annually across dozens of countries~\citep{accessnow2026keepiton}.
Iran's 2025--26 shutdown is among the most extended and heavily reported
examples of the period, but public reporting on \emph{how} such a shutdown
was implemented -- by what technical mechanism, in what sequence, and with
what selectivity -- is necessarily limited to what is externally
observable. This paper reconstructs the technical execution of Iran's
2025--26 shutdown from two independent classes of Internet-scale
measurement: BGP routing visibility (RouteViews and RIS route collectors)
and active-probing/network-telescope outage signals (IODA), joined on a
shared timeline per Autonomous System.

The reconstruction is organized around four questions, each mapped to one
governing hypothesis and tested against the joined measurement:

\begin{description}
\item[H1 (mechanism).] Was the blackout executed by withdrawing routes from
  BGP (a control-plane action, visible as routes disappearing) or by
  filtering traffic while routes remained announced (a data-plane action,
  visible as reachability collapsing without route withdrawal)?
  Section~\ref{sec:h1}.
\item[H2 (topology).] Was disconnection executed at a small number of
  international-gateway ASes, or distributed across individual access
  networks? Section~\ref{sec:h2}.
\item[H3 (selectivity).] During the partial, tiered restoration that began
  in P4, did government, banking, and state-media networks regain
  international visibility earlier and more completely than consumer/mobile
  networks? Section~\ref{sec:h3}.
\item[H4 (rehearsal).] Was the 2026 onset executed faster or more uniformly
  than Iran's own prior shutdown/restoration events (November 2019, June
  2025)? Section~\ref{sec:h4}.
\end{description}

Every threshold, exclusion rule, and definition this measurement depends on
-- what counts as ``visible,'' what counts as ``dark,'' where phase
boundaries fall, which population of ASes is in scope -- is decided and
recorded before its effect on any finding is seen, and is reported in full
in Section~\ref{sec:methods} together with the robustness checks each
decision obligates. Every empirical claim in Sections~\ref{sec:h1}--\ref{sec:h4}
carries an explicit attribution tier (observed / strongly implied /
consistent with reporting), and disconfirming evidence -- results that
complicate a hypothesis instead of confirming it -- is reported alongside
confirming evidence, not glossed over
(Section~\ref{sec:discussion}).

This study's headline finding -- a filtering-dominant blackout
(Section~\ref{sec:h1}) -- is also independently reached, via an entirely
different measurement stack, by a contemporaneous study of the same
event~\citep{jahromi2026multiperspective}. Part of the period this paper
covers -- an 18-day event in January 2026, identified in this study's own
data via a blind scan and not assumed from prior reporting -- is also
independently covered by a second
contemporaneous study~\citep{aceto2026iran}; this paper's contribution
regarding that specific event is an independent reproduction with finer
BGP-level timing, not a discovery claim.
